\documentclass[11pt]{article}

\usepackage[margin=1in]{geometry}
\usepackage{amsmath,amssymb}
\usepackage{booktabs}
\usepackage{graphicx}
\usepackage{microtype}
\usepackage{hyperref}
\usepackage{enumitem}
\hypersetup{colorlinks=true,allcolors=blue}
\setlist{nosep}
\graphicspath{{paper/figures/}{figures/}}

\title{Short Addition Completions in GPT-2 Small: Behavioral Limits and a Target-Dependent Equal-Operand Pattern}
\author{Anay Katiyar\\\small Independent Researcher}
\date{}

\begin{document}
\maketitle

\begin{abstract}
I examine short decimal-addition prompts in GPT-2 Small (124M parameters) using behavioral tests, attribution, interventions, and prompt comparisons. On a held-out zero-shot set, the model answers 2/32 prompts correctly; its best constant guess answers 8/32. In a 79-cell prompt grid, target parity matches the sign of the logit contrast in 68/79 cells, with agreement varying by target range. Across seven even sums, the mean equal-operand advantage is $+0.2474$ (two-sided exact sign-flip $p=0.03125$; target-level SE $=0.102$), while sign and control-exclusion checks give $p=0.125$. In a 20-target extension, the mean arithmetic contrast is $+0.39392$ (exact one-sided sign-flip $p=0.013618$; stratified 95\% interval $[+0.09556,+0.71692]$), but only 8/20 contrasts are positive, below the 70\% breadth criterion. The analysis also uses different target-selection and prompt/control details from those specified in the written plan. A post hoc split shows a larger mean among six multiples of ten; this pattern motivates further tests of target roundness and frequency, but establishes neither explanation nor an addition circuit.
\end{abstract}

\section{Introduction}

Mechanistic interpretability asks how a model's internal representations and computations produce its behavior. It aims to explain more than input-output correlations alone \cite{elhage2021framework}. GPT-2 is an autoregressive language model trained to predict the next token \cite{radford2019gpt2}. Arithmetic offers a clear test case because the intended operation is explicit. But a model's preference for the correct answer token does not, by itself, show that the model performed the operation.

Hanna et al. provide the closest comparison by model: they study greater-than prediction in pretrained GPT-2 Small \cite{hanna2023greater}. They identify and test a circuit for comparing year-like values, not for decimal addition. Other studies examine arithmetic mechanisms directly in pretrained language models. Stolfo et al. use causal mediation to trace arithmetic-related information flow in models with 2.8B--7B parameters \cite{stolfo2023causal}. Yu and Ananiadou analyze addition and other operations in Llama-7B and GPT-J, including attention heads and feed-forward neurons \cite{yu2024cna}. Nikankin et al. characterize arithmetic circuits and heuristic neurons in pretrained Llama 3, Pythia, and GPT-J models \cite{nikankin2025heuristics}. This prior work shows that researchers have already studied arithmetic mechanisms in pretrained models. Its models, scales, prompts, and analyses differ from my short decimal-completion probes in GPT-2 Small. Those circuit findings therefore neither validate nor rule out my interpretations. Nanda et al. and Zhong et al. offer another comparison: they study transformers trained on modular addition, where researchers can examine learned algorithms in controlled settings \cite{nanda2023progress,zhong2023clock}. Ye et al. study models trained on synthetic grade-school math, a distinct setting \cite{ye2024physics}.

I test GPT-2 Small with prompts such as \texttt{3 + 5 =}. I began with the hypothesis that addition might rely on a general, addition-specific circuit. Behavioral baselines, corrected attributions, and intervention controls narrowed that interpretation. The seven-target scan still showed a target-dependent preference for equal-operand prompts. I describe this as a behavioral effect and a candidate for future causal testing.

My contribution is narrower than a new account of arithmetic mechanisms. I document exact-answer baselines, a neighbor-token logit contrast, and an equal-operand comparison in this GPT-2 Small prompt setting. I also show how attribution and intervention audits limit circuit-level interpretations. The equal-operand pattern remains exploratory and descriptive. I do not claim priority for mechanistic arithmetic research, identify an addition algorithm, or conclude that GPT-2 Small lacks arithmetic-related representations. Nor is this focused comparison a systematic review of all arithmetic-interpretability work.

\section{Study scope and related work}

\paragraph{Pretrained mathematical behavior in GPT-2 Small.} Hanna et al. study greater-than prediction in pretrained GPT-2 Small and validate a circuit across related contexts \cite{hanna2023greater}. Their work is my closest precedent by model family, but not by arithmetic operation. I study decimal-addition prompts and do not provide comparable circuit validation or evidence of generalization.

\paragraph{Mechanistic studies of pretrained arithmetic.} Researchers have directly investigated arithmetic computations in pretrained models. Stolfo et al. use causal mediation on arithmetic questions. They report information transfer from operands and operators toward the answer position, followed by result-related processing in MLP modules \cite{stolfo2023causal}. Yu and Ananiadou study arithmetic heads and feed-forward neurons in Llama-7B, with related analyses in GPT-J \cite{yu2024cna}. Nikankin et al. analyze arithmetic circuits in pretrained Llama 3, Pythia, and GPT-J models. They also characterize groups of heuristic neurons that contribute to answer logits \cite{nikankin2025heuristics}. These studies directly inform mechanistic arithmetic analysis and the broader question of algorithmic computation versus narrower patterns. I do not replicate those studies. I use GPT-2 Small, short next-token decimal prompts, a low-accuracy behavioral setting, and an equal-operand contrast. My narrower contribution is to document that contrast and the limits of the evidence available to interpret it. Differences in models and methods mean that neither set of findings confirms the other.

\paragraph{Arithmetic in trained models and other tasks.} Nanda et al. study small transformers trained on modular addition. They use activation, weight, and intervention analyses to validate a learned Fourier algorithm \cite{nanda2023progress}. Zhong et al. show that architecture and training conditions shape the solutions that models learn for modular addition \cite{zhong2023clock}. Lee et al. examine how training data and formats affect arithmetic learning in small transformers \cite{lee2024teaching}. Ye et al. analyze internal processes in synthetic grade-school math settings, including GPT-2-family models trained on generated data \cite{ye2024physics}. These studies offer useful comparisons, but their training regimes or tasks differ from my pretrained short-completion setting.

\paragraph{Circuit-evaluation methods.} Studies of indirect-object identification and automated circuit discovery motivate explicit behavioral metrics and causal validation \cite{wang2023ioi,conmy2023acdc}. They provide methodological context. All measurements below come from the project data.

\section{Methods}

\subsection{Model, prompts, and data}

I study GPT-2 Small (124M parameters) using short arithmetic prompts, mainly \texttt{a + b =}. The repository contains my full-record notebook, analysis outputs, raw corpus query logs, a 79-row addition grid, a 39-row digit-addition double/control score file, and a checksummed artifact manifest. The format and operator tables contain target-level summaries, but not the underlying prompt-level scores.

I report three behavioral cohorts separately: all-pairs zero-shot ($N=36$), held-out zero-shot ($N=32$), and held-out few-shot ($N=32$). I do not pool them. “Held-out” follows the cohort definition used here: it excludes the four ordered operand pairs used as few-shot examples and their order-reversed counterpart, $(1,1)$, $(2,2)$, $(5,3)$, and $(3,5)$. The few-shot prefix contains three demonstrations, each on its own line: \texttt{1 + 1 = 2}, \texttt{2 + 2 = 4}, and \texttt{5 + 3 = 8}. The all-pairs cohort contains every ordered single-digit pair with a one-digit sum. Under the exclusions above, each held-out cohort contains the remaining 32 pairs. The dense equal-operand scan covers seven even target sums, $T\in\{4,6,8,10,12,14,16\}$. For each target, I compare one double prompt, $(T/2)+(T/2)$, with the available ordered unequal single-digit splits that sum to $T$. The exploratory two-digit transfer analysis covers targets 24, 26, 28, 30, 32, 34, and 36. For each target, it compares one double prompt with ordered unequal splits whose operands range from 10 to 19. I froze the three-component candidate set before this transfer analysis. Because earlier analyses selected these components, the transfer analysis is not an independent confirmation.

I extend the target scan to 20 even sums from 38 to 98: 38, 40, 46, 48, 50, 54, 56, 58, 60, 68, 70, 72, 74, 78, 80, 84, 86, 90, 94, and 98. I use leading-space target and numeric-neighbor tokens and the neighbor-averaged score in Eq.~\ref{eq:score}. Each target has one arithmetic prompt template and one unequal split, evaluated in both operand orders. The saved records also include one comma template and four alternate natural-language templates; I treat those formats as exploratory. The selection used seed 42 and ranges 38--56, 58--76, and 78--98, rather than the written plan's seed 2026101201 and ranges 38--58, 60--78, and 80--98. Under the plan's ranges, the selected targets fall into bins of 8/6/6, rather than 7/6/7. The analysis also includes fewer arithmetic templates and unequal splits than planned. I report it as a continuation of the same study while keeping these design differences visible.



\subsection{Behavioral metrics}

I measure behavioral performance by full-vocabulary top-1 accuracy and compare it with a cohort-specific constant baseline. For prompt $p$ and target token $T$, the neighbor-averaged logit contrast is
\begin{equation}
S(p,T)=\ell_T(p)-\frac{\ell_{T-1}(p)+\ell_{T+1}(p)}{2},
\label{eq:score}
\end{equation}
where $\ell_j(p)$ is the next-token logit for the leading-space answer token corresponding to $j$. This score measures relative logit preference among the target and its two neighboring values; it is not a measure of exact-answer probability or proof of arithmetic computation.

For each even target $T$, the equal-operand advantage is
\begin{equation}
A_T=S(p_{d,d},T)-\frac{1}{|C_T|}\sum_{(a,b)\in C_T}S(p_{a,b},T),\qquad d=T/2,
\label{eq:adv}
\end{equation}
Here, $p_{a,b}$ is the prompt \texttt{a + b =}, and $C_T$ contains the evaluated ordered unequal splits whose operands sum to $T$. I summarize the seven target-level advantages with their mean and standard error; the target is the inferential unit ($n=7$). I use a two-sided exact sign-flip test over all $2^7$ sign assignments and a 10,000-resample percentile bootstrap over targets (seed 20261008). With only seven target clusters, the bootstrap interval is necessarily coarse.

For the causal screen, I used five discovery targets, $T\in\{4,6,10,12,16\}$, and reserved targets 8 and 14 for exploratory transfer. I selected 12 components with the largest absolute DLA contrasts for nominated-component mean-ablation screening. For each component, I measured how ablation changed the equal-operand advantage and used a one-sided exact sign-flip test to ask whether ablation reduced it. I applied Benjamini--Hochberg correction across all 12 components. Four components reached the smallest possible one-sided $p=1/32=0.03125$. I selected the three with the largest negative mean ablation changes (\texttt{10\_mlp\_out}, L9H9, and L10H2) for exploratory transfer. With five target-level units, p-value resolution is coarse: the smallest two-sided exact sign-flip p-value is $2/32=0.0625$. In a 12-test BH family, at least 8 tests must reach the one-sided floor for $q<0.05$; two-sided $q<0.05$ is unattainable. This screen therefore cannot confirm a candidate. Failure to pass correction also does not show that no relevant component exists.

For the addition grid, I record whether the measured symmetric logit-contrast sign matches target parity. Grid cells share target values, so they are not independent observations. In a separate sensitivity analysis, I replace foils $T-1,T+1$ with $T-2,T+2$ and report the value at each target. I define the component-completeness ratio at each discovery target as $(A_T^{\mathrm{full}}-A_T^{\mathrm{ablated}})/A_T^{\mathrm{full}}$. This ratio can become unstable when the full-model contrast is small or changes sign.

\subsection{Attribution and intervention audits}

In my attribution analysis, I project residual components onto a target-versus-foil unembedding direction. An audit corrected an earlier calculation: it applied final LayerNorm scaling to the residual decomposition before projection and separately included the unembedding bias when reconstructing the measured logit contrast. These corrections change how I interpret selected prompts. Attribution alone does not establish necessity or sufficiency.

I compare zero ablation with mean ablation in the intervention summaries. Candidate screening uses mean ablation and reports false-discovery-rate-adjusted $q$ values. In the exploratory transfer analysis, I apply the three frozen components to double and matched unequal-operand prompts at each target.

I checked filler specificity with four fixed non-arithmetic prompts: \texttt{cat dog bird =}, \texttt{The object in the box =}, \texttt{The word on the page =}, and \texttt{Yesterday at the store =}. A separate operator-position patching experiment lacked a corrupt-baseline anchor. It therefore cannot establish that operator binding was falsified. I report that experiment separately from the aggregate operator-string comparisons.

\subsection{Format, operator, and corpus comparisons}

I archive seven matched target-level means per condition for the format and operator analyses, along with paired target-level tests and Holm correction. A separate notebook summary reports pooled control standard deviations over four targets for three numeral formats. These dispersions and their advantage/SD ratios are descriptive, not significance tests. The operator and connector conditions use the source strings \texttt{+}, Unicode subtraction (\texttt{−}) and multiplication (\texttt{×}), \texttt{and}, and \texttt{then}. For strings other than addition, I still score the arithmetic sum token. These comparisons do not test whether the model performs those operations.

I use Infini-gram counts from the Dolma v1.7 index \texttt{v4\_dolma-v1\_7\_llama} as a proxy for textual co-occurrence. The query set includes double strings and matched unequal-operand controls; I apply no minimum-count threshold. Afterward, I computed an exploratory Spearman analysis from the seven-target summary and the equal-operand advantages. I had not prespecified this analysis. I archive the raw request/response records and the summary without a count threshold. The API tokenization record separates the first operand digit from its preceding space. As a result, exact-string counts may also include longer numbers that end in the queried sequence. Dolma is not GPT-2's exact training corpus.

\section{Results}

\subsection{The tested cohorts do not show reliable exact-answer performance}

The model achieves full-vocabulary top-1 accuracy of 2/36 (5.56\%) on the all-pairs zero-shot cohort and 2/32 (6.25\%) on the held-out zero-shot cohort. The best constant guess scores 8/36 (22.2\%) and 8/32 (25\%), respectively. Held-out few-shot full-vocabulary accuracy is also 2/32. Restricting prediction to the seven answer digits in the held-out cohort raises top-1 accuracy to 4/32 (12.5\%) in zero-shot and 10/32 (31.25\%) in few-shot. A constant digit guess scores 8/32 (25\%) in either setting. The recorded one-sided comparison does not find a significant few-shot difference under this restricted-vocabulary metric ($p\approx0.26$).

A separate label-permutation analysis reports a row-centred own-answer logit contrast of $+0.224$ zero-shot and $+0.521$ few-shot (Monte Carlo $p=0.0004$ in both settings; 5,000 permutations). These restricted-vocabulary and logit-contrast results show answer-token signal under their respective metrics. They do not measure full-vocabulary exact-answer accuracy. My cohort-specific results do not show that GPT-2 Small lacks all arithmetic-related computation. They show that these prompts do not support reliable exact-answer addition above the simple baseline.

\subsection{Parity agreement varies across the addition grid}

In the committed 79-cell grid, target parity predicts the measured logit-contrast sign in 68/79 cells. Agreement is 63/64 for sums $\leq 13$ and 5/15 for higher-sum rows (sums $\geq 14$). Cells share target sums and differ in prompt composition, so these fractions are descriptive; they are not estimates from independent trials. The pattern also varies across the tested range. I ran a separate sensitivity analysis with parity-matched foils, $T-2$ and $T+2$:

\begin{center}
\small
\begin{tabular}{rr}
\toprule
Target $T$ & Advantage with parity-matched foils \\
\midrule
4 & $-0.1112$ \\
6 & $-0.1793$ \\
8 & $+0.1487$ \\
10 & $+0.4641$ \\
12 & $+0.1597$ \\
14 & $+0.1868$ \\
16 & $+0.6369$ \\
\bottomrule
\end{tabular}
\end{center}

The mean is $+0.1865$. The two-sided exact sign-flip test gives $p=0.140625$. The 95\% bootstrap interval is $[-0.0068,\allowbreak+0.3863]$. Under this foil definition, targets 4 and 6 are negative, while target 14 is positive. This sensitivity test does not reject a zero mean.

\begin{figure}[t]
\centering
\includegraphics[width=0.88\linewidth]{addition_grid_heatmap.pdf}
\caption{Addition-grid symmetric logit differences across the 79 prompt pairs. The grid shows target-dependent structure; cells sharing a target are not independent observations. Rendered from the canonical grid data.}
\label{fig:addition-grid}
\end{figure}

\subsection{Attribution corrections narrow component interpretations}

For the recorded \texttt{3 + 5 =} comparison, corrected residual-component DLA sums to $+0.4658$. The reported unembedding-bias contrast is $+0.1782$. Together, they reconstruct $+0.6440$, compared with a measured logit difference of $+0.6443$. I retain the small discrepancy as reported. In a separate few-shot comparison on \texttt{158 + 274 =}, the dynamic contribution is $+0.8382$ and the bias term is $-1.0303$, giving a reported logit difference of $-0.1920$. This contrast compares the leading-space tokens \texttt{8} and \texttt{1}, even though the arithmetic answer is 432. It is an identity between token contrasts, not evidence that the model computed 8 and was then overridden.

Zero ablation changed the residual-stream scale in the reported prompt. Across reference prompts, mean ablation gives a reported L11H0 shift of $+0.0013$. Jointly ablating L11H0 and MLP10 gives $+0.8127$, compared with the additive prediction of $+0.8131$ for that prompt. These results do not support the earlier suppressor and non-additivity interpretations for this prompt and intervention.

\subsection{A target-dependent equal-operand preference remains}

The 39-row digit--digit scan reports the following target-level summaries:

\begin{table}[t]
\centering
\caption{Target-level equal-operand advantage for the seven even sums. Exact sign-flip inference uses seven target-level contrasts.}
\label{tab:doubles-targets}
\small
\begin{tabular}{rrrrr}
\toprule
Target $T$ & Double score & Control mean & Advantage $A_T$ & Controls \\
\midrule
4  & $+0.251$ & $+0.203$ & $+0.049$ & 2 \\
6  & $+0.230$ & $+0.213$ & $+0.017$ & 4 \\
8  & $+0.606$ & $+0.315$ & $+0.291$ & 6 \\
10 & $+0.647$ & $+0.113$ & $+0.535$ & 8 \\
12 & $+0.401$ & $+0.236$ & $+0.165$ & 6 \\
14 & $-0.238$ & $-0.232$ & $-0.006$ & 4 \\
16 & $+0.719$ & $+0.037$ & $+0.682$ & 2 \\
\bottomrule
\end{tabular}
\end{table}

Across seven target levels, the mean advantage is $+0.2474$, and six target-level values are positive. The exact target-level sign-flip test gives $p=0.03125$. A 10,000-resample target bootstrap gives a 95\% percentile interval of $[+0.0738,+0.4374]$. The small-sample result is sensitive to the test and controls. The exact sign test for six positive targets out of seven gives $p=0.125$, while leave-one-target-out sign-flip tests give $p=0.03125$ or $0.0625$. I also excluded unequal prompts that use $T-1$ or $T+1$ as an operand. Six targets then remain because target 4 has no eligible controls. Their mean advantage is $+0.2561$, with exact sign-flip $p=0.125$. I therefore treat this as an exploratory pattern, not a robust effect across target and control specifications.

Across the seven target levels, the paired digit+digit versus digit+word difference is $-0.2190$ (raw $p=0.015625$; Holm-adjusted $p=0.046875$ within the three-comparison format family). This adjusted value is just below 0.05, and I did not adjust it across the wider set of project analyses. The other two format comparisons have Holm-adjusted $p=0.21875$. A separate four-target notebook summary reports pooled control SDs of $0.080$, $0.171$, and $0.273$ for digit+digit, digit+word, and word+word, respectively. It reports descriptive advantage/SD ratios of $5.22$, $4.40$, and $4.93$. Both raw advantages and control spread increase with word-form use. These ratios therefore do not establish a format effect, and they are not significance tests.

I have seven target-level means for each of five operator and connector strings. At target 8, the advantages for \texttt{and} ($+0.321$) and \texttt{then} ($+0.344$) are at least as large as the plus advantage ($+0.291$). This exploratory pattern is consistent with an effect that is not addition-specific. However, all four plus-versus-other paired tests have Holm-adjusted $p=0.6875$. I detect no difference at this sample size, but that does not establish equivalence. The task also lacks semantic matching: I score the same arithmetic sum token under addition, subtraction, multiplication, and connectives. These comparisons therefore do not show that the model performs arithmetic under those strings. The prompt-level operator scores are not included.

\begin{figure}[t]
\centering
\includegraphics[width=0.96\linewidth]{operator_target_advantages.pdf}
\caption{Target-level mean equal-operand advantage for each operator/connector string across seven targets. Same sum token is scored in every condition; this does not test whether the model computes each operation.}
\label{fig:operator-summary}
\end{figure}

\subsection{Twenty-target extension: positive mean but limited breadth}

For this 20-target analysis, I use spacing-aware target and numeric-neighbor tokens. The per-target margins match the corresponding target and neighbor logits. The mean target-level arithmetic contrast is $+0.39392$. Under the written plan's original ranges, the sample allocation is 8/6/6. The 95\% bootstrap interval is $[+0.09556,+0.71692]$, and the exact one-sided sign-flip test gives $p=0.013618$ (two-sided $p=0.027235$). The mean exceeds the $0.10$ threshold, and the interval and test criteria pass. Only 8/20 targets ($40\%$) have positive contrasts, below the 70\% threshold (14/20). Thus, three numerical criteria pass and the breadth criterion fails. The selection and prompt/control design differ from the written plan, so this analysis does not satisfy every planned condition.

In a post hoc analysis, the arithmetic contrast is concentrated in six selected targets divisible by ten (40, 50, 60, 70, 80, and 90). Their mean is $+1.4313$, and all six are positive. The other 14 targets have a mean of $-0.0507$, with two positive values; the overall median is $-0.0606$. This contrast largely tracks the equal-prompt scores. Their mean is $+1.7410$ for round targets and $+0.0022$ for other targets. The corresponding unequal-control means are $+0.3097$ and $+0.0529$. A fixed-size, two-sided permutation of the six round labels gives an unadjusted $p=1/38{,}760\approx2.6\times10^{-5}$ for the arithmetic-contrast gap. I identified this subgroup after inspecting the results, and I did not adjust for searching across possible subgroup definitions. I therefore treat it as exploratory, not confirmatory. The same descriptive round/non-round split appears in the comma and four alternate natural-language formats. These measurements come from the same targets and are correlated, so they are not independent replications. Across targets, the arithmetic-minus-comma contrast averages $-0.134$ (correlation $r=0.878$). A nonsignificant comparison does not establish equivalence. These prompt formats also differ from those specified in the written plan.

My notebook also reports a summary for the 11 remaining targets: a pooled-across-formats mean of $+0.282$, one-sided sign-flip $p=0.175$, and 7/11 positive values. None of these targets is a multiple of ten. This analysis pools six formats instead of isolating the arithmetic condition. The repository also lacks its prompt-level output file. I therefore treat it neither as an independent confirmation nor as a test of the round-target pattern.

\subsection{Candidate-component screens do not establish a circuit}

My notebook screens 12 components, selected by the magnitude of their DLA contrasts across five discovery targets. Four components reach the smallest possible one-sided p-value, $0.03125$. I selected the three with the largest negative mean changes in advantage after ablation—\texttt{10\_mlp\_out}, L9H9, and L10H2—for subsequent exploratory tests. L8H11 is the fourth component at the p-value floor; I did not select it because its mean change ranked fourth. For the selected three, the one-sided BH-adjusted $q=0.09375$ and two-sided $q=0.15$. None passes $q<0.05$. With five target-level units and a 12-test family, this screen cannot confirm a component. Its failure to pass correction also does not count as evidence against one. Target-level completeness ratios range from $0.54$ to $1.54$ and become unstable when the full-model contrast is small.

In the exploratory transfer analysis at two-digit targets, the clean and jointly ablated equal-operand advantages were:

\begin{center}
\small
\begin{tabular}{rrrr}
\toprule
Target & Clean advantage & Ablated advantage & Change \\
\midrule
24 & $+0.446$ & $+0.239$ & $-0.207$ \\
26 & $-0.117$ & $-0.026$ & $+0.091$ \\
28 & $+0.268$ & $+0.107$ & $-0.161$ \\
30 & $+0.770$ & $+0.472$ & $-0.298$ \\
32 & $+0.767$ & $+0.413$ & $-0.354$ \\
34 & $-0.066$ & $+0.035$ & $+0.101$ \\
36 & $+0.624$ & $+0.309$ & $-0.314$ \\
\bottomrule
\end{tabular}
\end{center}

The clean advantage is positive at five of seven targets (mean $+0.384$; two-sided exact sign-flip $p=0.0625$). The three-component ablation reduces the advantage at those five targets and increases it at targets 26 and 34. This exploratory result extends beyond the discovery range, but it is not an independent confirmation: earlier analyses selected the components. Filler prompts also show component effects, so I have not established specificity. Taken together, localization, ablation, patching, transfer, and filler checks do not establish necessity, sufficiency, minimality, or a validated addition circuit.

\subsection{Corpus-frequency relationship remains unresolved}

My corpus query log records 79 API attempts. All returned HTTP 200, and none has a logged error. The archived summary, which applies no count threshold, shows positive double-versus-control joint log ratios at all seven targets and mixed conditional log ratios. In a post hoc Spearman analysis across these targets, joint log ratio versus equal-operand advantage gives $\rho=+0.464$ ($p=0.294$); conditional log ratio versus advantage gives $\rho=-0.143$ ($p=0.760$). These small-sample tests detect no association. I did not prespecify them, and they do not rule out a frequency contribution. The counts also have an exact-string digit-boundary limitation: the API tokenizer records the first digit without a leading-space token marker, so a match may include a longer number ending in the queried digits. Dolma is a proxy, not GPT-2's training corpus.

\section{Discussion}

I began by asking whether an addition circuit could explain the observed behavior. The tested cohorts do not show reliable exact-answer performance, and the parity grid shows that logit preferences vary across target ranges. Attribution and intervention audits narrowed several component-level explanations. The 20-target extension has a positive mean arithmetic contrast, but only 8/20 targets are positive, so it fails the breadth threshold. A post hoc subgroup of multiples of ten accounts for much of the contrast, and alternate prompt formats show similar patterns. This points to target-dependent behavior, but it does not establish a frequency mechanism or explain the source of the logits. The analysis also differs from the written plan in target selection and prompt/control design, which limits its confirmatory interpretation.

I interpret these findings alongside prior mechanistic studies of arithmetic in pretrained models, including causal mediation, attention/MLP analyses, and heuristic-neuron accounts \cite{stolfo2023causal,yu2024cna,nikankin2025heuristics}. My measurements neither reproduce nor challenge those studies because they use different models and prompt/evaluation setups. The closest comparison by model family is the GPT-2 Small greater-than study, which validates a circuit for a different task \cite{hanna2023greater}. Studies of modular addition offer a complementary comparison using controlled training \cite{nanda2023progress,zhong2023clock}. I therefore present a model- and prompt-specific behavioral audit with an exploratory equal-operand contrast, not a general theory of how language models perform arithmetic.

The most direct follow-up would test the post hoc round-number pattern in a new sample fixed in advance. It should balance round and non-round sums, include several unequal splits per target, and use all planned arithmetic templates. It should vary operand roundness and target roundness separately. I should retain prompt-level outputs, apply prespecified selection rules, and compare arithmetic with matched repetition controls using target-level inference. Causal interventions can follow; the present results do not identify a mechanism.

\section{Limitations}

I examine one model, a narrow set of short prompts, and a small number of target sums. The cohort definitions and few-shot prefixes differ, so I report them separately. The repository includes raw corpus API responses. However, the format and operator exports contain target-level summaries rather than full prompt-level score tables. The corpus association analysis is post hoc, covers only seven targets, and relies on a Dolma proxy with a digit-boundary limitation. I have not tested the single-operand comparison. The candidate-component results are exploratory and cannot support a circuit claim.

The 20-target extension differs from the written plan: it used another selection seed and target-bin boundaries, one arithmetic template instead of three, and one unequal split in both orders instead of up to six sampled splits. Under the plan's original bins, the targets divide 8/6/6, while the stored 7/6/7 labels use revised ranges. The comma and natural-language templates also differ from the planned controls. The score file checks leading-space target and neighbor tokens, but it does not preserve the full context-by-context operand-tokenization audit specified in the plan. The analysis fails the 70\% positive-target criterion. Its round-number subgroup is post hoc and includes only six targets; equal splits also confound target roundness with operand roundness. I did not query corpus counts for these targets. The 11-target analysis comes from an aggregate notebook output and does not isolate the arithmetic template. I tested no other model. In addition, the operator comparisons score the arithmetic-sum token even when a prompt uses a different operation, and the corpus proxy is not GPT-2's training corpus. These limits constrain mechanistic and generalization claims.

\section{Conclusion}

For the prompts I evaluated, the evidence does not establish a general addition-specific circuit in GPT-2 Small. The 20-target extension shows a positive mean arithmetic contrast, but only 8/20 target contrasts are positive, and the analysis uses fewer templates and controls than the written plan specified. A post hoc split shows a larger mean for multiples of ten and a similar descriptive pattern in alternate prompt formats; it does not explain why the pattern occurs. These findings establish neither a broad equal-operand effect nor a specific frequency or arithmetic mechanism. They motivate a controlled test that separates target roundness from operand roundness.

\section*{Data and code availability}

My research notebook, score tables, audit trail, and scripts are available at \url{https://github.com/anaykatiyar-lang/mechanistic_interpretability}. The repository contains analysis outputs, raw corpus request/response logs, and a SHA-256 data manifest. The format and operator artifacts provide target-level summaries; the manifest describes each output.

\section*{AI assistance disclosure}

I used Claude Sonnet 5.5, ChatGPT, Gemini 3.1 Pro, and Codex to explore literature, review and generate code, and edit the manuscript. I treated their suggestions as prompts for review, not as evidence. AI-assisted drafts contained bibliographic inaccuracies; I checked the cited records against primary sources and corrected them during review. I checked reported quantities against my analysis outputs and take responsibility for the experiments, data, analysis, claims, and final manuscript.

\begin{thebibliography}{99}

\bibitem{conmy2023acdc}
A.~Conmy, A.~N. Mavor-Parker, A.~Lynch, S.~Heimersheim, and A.~Garriga-Alonso.
\newblock Towards automated circuit discovery for mechanistic interpretability.
\newblock In \emph{Advances in Neural Information Processing Systems}, 2023.
\newblock \url{https://arxiv.org/abs/2304.14997}.

\bibitem{elhage2021framework}
N.~Elhage et~al.
\newblock A mathematical framework for transformer circuits.
\newblock \emph{Transformer Circuits Thread}, 2021.
\newblock \url{https://transformer-circuits.pub/2021/framework/index.html}.

\bibitem{hanna2023greater}
M.~Hanna, O.~Liu, and A.~Variengien.
\newblock How does GPT-2 compute greater-than? Interpreting mathematical abilities in a pre-trained language model.
\newblock In \emph{Advances in Neural Information Processing Systems}, 2023.
\newblock \url{https://arxiv.org/abs/2305.00586}.

\bibitem{lee2024teaching}
N.~Lee, K.~Sreenivasan, J.~D. Lee, K.~Lee, and D.~Papailiopoulos.
\newblock Teaching arithmetic to small transformers.
\newblock In \emph{International Conference on Learning Representations}, 2024.
\newblock \url{https://arxiv.org/abs/2307.03381}.

\bibitem{nanda2023progress}
N.~Nanda, L.~Chan, T.~Lieberum, J.~Smith, and J.~Steinhardt.
\newblock Progress measures for grokking via mechanistic interpretability.
\newblock In \emph{International Conference on Learning Representations}, 2023.
\newblock \url{https://arxiv.org/abs/2301.05217}.

\bibitem{nikankin2025heuristics}
Y.~Nikankin, A.~Reusch, A.~Mueller, and Y.~Belinkov.
\newblock Arithmetic without algorithms: Language models solve math with a bag of heuristics.
\newblock In \emph{International Conference on Learning Representations}, 2025.
\newblock \url{https://arxiv.org/abs/2410.21272}.

\bibitem{radford2019gpt2}
A.~Radford et~al.
\newblock Language models are unsupervised multitask learners.
\newblock OpenAI technical report, 2019.
\newblock \url{https://cdn.openai.com/better-language-models/language_models_are_unsupervised_multitask_learners.pdf}.

\bibitem{stolfo2023causal}
A.~Stolfo, Y.~Belinkov, and M.~Sachan.
\newblock A mechanistic interpretation of arithmetic reasoning in language models using causal mediation analysis.
\newblock In \emph{Proceedings of the 2023 Conference on Empirical Methods in Natural Language Processing}, pages 7035--7052, 2023.
\newblock \url{https://aclanthology.org/2023.emnlp-main.435/}.

\bibitem{wang2023ioi}
K.~Wang, A.~Variengien, A.~Conmy, B.~Shlegeris, and J.~Steinhardt.
\newblock Interpretability in the wild: A circuit for indirect object identification in GPT-2 small.
\newblock In \emph{International Conference on Learning Representations}, 2023.
\newblock \url{https://arxiv.org/abs/2211.00593}.

\bibitem{yu2024cna}
Z.~Yu and S.~Ananiadou.
\newblock Interpreting arithmetic mechanism in large language models through comparative neuron analysis.
\newblock In \emph{Proceedings of the 2024 Conference on Empirical Methods in Natural Language Processing}, pages 3293--3306, 2024.
\newblock \url{https://aclanthology.org/2024.emnlp-main.193/}.

\bibitem{ye2024physics}
T.~Ye, Z.~Xu, Y.~Li, and Z.~Allen-Zhu.
\newblock Physics of language models: Part 2.1, grade-school math and the hidden reasoning process.
\newblock arXiv:2407.20311, 2024.
\newblock \url{https://arxiv.org/abs/2407.20311}.

\bibitem{zhong2023clock}
Z.~Zhong, Z.~Liu, M.~Tegmark, and J.~Andreas.
\newblock The clock and the pizza: Two stories in mechanistic explanation of neural networks.
\newblock In \emph{Advances in Neural Information Processing Systems}, 2023.
\newblock \url{https://arxiv.org/abs/2306.17844}.

\end{thebibliography}
\end{document}
